\chapter{应用落地与维护策略}

\section{部署模式}
推荐“两阶段部署”：
\begin{itemize}
  \item 离线阶段：建立标称参数与可行域；
  \item 在线阶段：基于滚动窗口进行轻量修正。
\end{itemize}

\section{监控指标}
在线系统应持续监控：
\begin{enumerate}
  \item 残差均值和方差；
  \item 参数漂移速度；
  \item 迭代失败率与重启次数。
\end{enumerate}
这些指标可用于判断当前模型是否失配。

\section{维护机制}
建议按周执行全量回放校验，按月更新先验参数范围。对长期漂移场景，应引入缓慢时变参数模型。

\section{面向工程的结论}
模型层、算法层和系统层必须协同优化：仅提升单点精度不足以保障长期稳定运行，需同时关注鲁棒性、可解释性与可维护性。
